% 原书真实矢量数据及标记字形；中文只覆盖数据区外的标签。
\begin{tikzpicture}[x=0.51mm,y=-0.51mm]
\node[anchor=north west,inner sep=0pt] at (0,0) {\includegraphics[width=150.9600mm,height=98.9400mm]{figures/ch28_fig01_data.pdf}};
\node[font=\fontsize{9.5}{10.5}\selectfont,inner sep=0pt,anchor=east] at (88.24687,155.37414) {联合靶向自助法};
\node[font=\fontsize{9.5}{10.5}\selectfont,inner sep=0pt,anchor=east] at (88.24831,128.15646) {靶向自助法};
\node[font=\fontsize{9.5}{10.5}\selectfont,inner sep=0pt,anchor=east] at (88.24117,100.94662) {超级学习自助法};
\node[font=\fontsize{9.5}{10.5}\selectfont,inner sep=0pt,anchor=east] at (88.24329,73.72897) {非参数自助法 + CVID};
\node[font=\fontsize{9.5}{10.5}\selectfont,inner sep=0pt,anchor=east] at (88.24332,46.51773) {非参数自助法};
\node[font=\fontsize{9.5}{10.5}\selectfont,inner sep=0pt,anchor=east] at (88.24402,19.30008) {完整样本};
\node[font=\fontsize{9.5}{10.5}\selectfont,inner sep=0pt,anchor=center] at (193.41423,187.65699) {风险差比值};
\node[font=\fontsize{9.5}{10.5}\selectfont,inner sep=0pt,anchor=center,rotate=90] at (6.86757,86.71398) {样本类型};
\end{tikzpicture}
